% Supplement subsection for study IJ (from ij_section.draft.tex; Reading written after the numbers, 2026-10-25).  To use it once experiments/results/ij_summary.json exists:
%   cp docs/paper/ij_section.draft.tex docs/paper/ij_section.tex  (then edit the "Reading" paragraph to what the numbers show).  The include is added to supplement.tex by \IfFileExists.
\subsection{Study IJ: An Answerability Probe for Unknown Destinations}\label{sec:moved-ij}
IG and IH showed that the parser reads requests to unknown places as requests to shelter~B, and that neither a list rule nor the model's own destination probability separates them. Probes on hidden states can detect unanswerable or out-of-scope inputs where output probabilities fail~\cite{wagner2026twoaxes,lavi2025linear}: models encode answerability even while answering an unanswerable question confidently~\cite{slobodkin2023unanswerability}, and hidden states beat sentence probabilities for truthfulness~\cite{azaria2023lying,marks2023geometry}. The same literature warns that probes trained on one dataset or task transfer poorly~\cite{orgad2025knowmore,azizian2025orthogonal} and track knowledge recall more than truth~\cite{cheang2025recall}, which is why IJ holds out whole name families rather than single names. Study IJ (registered and stamped before any feature was extracted) fits a logistic probe on the frozen Qwen3.5-9B's hidden states (layer \IJLayer{} chosen on development data by leave-one-family-out AUROC) to flag requests to unknown destinations; \IJDev{} development and \IJTest{} test texts, five whole place-name families held out for test. Thresholds come from the development negatives by the rank rule.
\emph{Claims.} IJ1: probe AUROC \IJProbeAuroc{} (\IJProbeAurocCI, family bootstrap) at least 0.05 above the better of the model's zero-shot yes-probability (\IJZeroShotAuroc) and its closed-set probability of ``other'' (\IJPOtherAuroc): \IJOneHolds. IJ2: at most 15\% of unknown-destination requests pass the probe (\IJProbePass): \IJTwoHolds. IJ3: at most 8\% of listed-destination requests refused (\IJProbeFalseRef): \IJThreeHolds. Baselines: difference of means \IJDiffMeansAuroc{} (pass-through \IJDiffMeansPass), list rule pass-through \IJListRulePass. Predictions: the probe passes fewer than the list rule \IJPredOneHolds; at least 4 of 5 test families pass through at most 25\% \IJPredTwoHolds.
\emph{Reading.} The hidden state knows what the output hides, but the threshold does not travel. On the held-out name families the probe separates unknown from listed destinations far better than the model's own probabilities (AUROC \IJProbeAuroc{} against \IJZeroShotAuroc{} and \IJPOtherAuroc; IJ1 \IJOneHolds) and refuses no listed destination (\IJProbeFalseRef; IJ3 \IJThreeHolds), yet \IJProbePass{} of the unknown destinations pass it (IJ2 \IJTwoHolds): the rank-rule threshold set on the development negatives sits above most test scores, because the development families were separated less well (leave-one-family-out AUROC on development data about 0.70). The difference of means and the list rule let fewer through (\IJDiffMeansPass, \IJListRulePass) at a lower AUROC. So a probe is a promising signal for out-of-list destinations, but a deployable threshold needs calibration data from the site's own place names, which is the transfer warning of the probing literature in a small, registered form. The parser stays as submitted and the read-back of supply and destination stays mandatory; we do not run a further study (experiments are closed).
